\documentclass[12pt]{article}
\usepackage{amsmath, amssymb, mathrsfs}
\usepackage[T1]{fontenc}
\usepackage[hmargin=1in,vmargin=1in]{geometry}
\usepackage{amsthm}
\newtheorem{theorem}{Theorem}[section]
\newtheorem{lemma}[theorem]{Lemma}
\usepackage{hyperref}

\begin{document}

\begin{center} {\large \textbf{Math 441A -- Homework 2}} \smallskip \\ 
    Grant Yang \hspace{1em} AU26
 \end{center}

\paragraph{Problem 1.} Consider Euclidean $n$-space $(\mathbb R^n, d)$, and let $P, Q, R \in \mathbb R^n$. If we have that \begin{equation}d(P, Q) + d(Q, R) = d(P, R), \label{p1} \end{equation} then $Q$ lies on the line segment $PR$.

\begin{proof}
    Let $|x|^2 = \langle x, x \rangle$ denote the square norm. We can calculate that
    \begin{align*}
        d(P, R) &= |P - R| = |(P - Q) - (R - Q)| \\
        &= \sqrt{|P - Q|^2 + |R - Q|^2 - 2 \langle P - Q, R - Q \rangle}.
    \end{align*}
    Then, we can apply the fact that $d(P, R) = d(P, Q) + d(Q, R) = |P - Q| + |Q - R|$ and square both sides:
    \begin{align*}
        |P-Q|^2 + |Q - R|^2 + 2|P-Q| |Q - R| = |P-Q|^2 + |R-Q|^2 + 2 \langle P - Q, Q - R \rangle.
    \end{align*}
    Thus, we must have that $|P - Q||Q - R| = \langle P - Q, Q - R \rangle$, implying that there exists some $\lambda$ such that $P - Q = \lambda (Q - R)$ as we have achieved equality in Cauchy Schwartz. Rearranging, we get \[Q = \frac{P + \lambda R}{\lambda + 1} = \frac{1}{\lambda + 1} P + \left(1 - \frac{1}{\lambda + 1}\right)R.\]
    Now, it suffices to show that $0 \le \frac{1}{\lambda + 1} \le 1$. We can compute \[|P - R| = |(P - Q) - (R - Q)| = |\lambda (Q - R) + (Q - R)| = |\lambda + 1| \cdot |Q - R|.\]
    From the positivity of the distance function and our equation (\ref{p1}), we know that \\ $|P - Q| \le |P - R|$. Thus, we get 
    \[ |\lambda| \cdot |Q - R| \le |\lambda + 1| \cdot |Q - R|, \]
    so $\lambda \ge -\frac{1}{2}$, and $0 \le \frac{1}{\lambda + 1} \le 1$ as desired.
\end{proof}

\paragraph{Problem 2.}


\paragraph{Problem 3.}
As subsets of $\mathbb R^2$ with the usual topology:
\begin{enumerate}
    \item[(1)] $A = \{(x, y) : x < 0\} \cup \{(x, y): x > 0, y > 1/x \}$ is open, as it is the union of two open sets. Let $(x, y) \in A$. If $x < 0$, then $B((x, y), |x|/2) \subseteq A$. Otherwise, $x > 0$, and you can find a ball too but I dont want to do the math.
    \item[(2)] Topologist's sine curve, $A = \{(x, \sin(1/x)) : x > 0\}  \cup \{(0, y) : -1 \le y \le 1\}$. I think this is closed?
    \item[(3)] $A = \{(x, y) : y \in \mathbb Q, \exists n \in \mathbb Z_{>0}, y = x^n\}$. This set is neither open nor closed. Clearly this thing is not open: note $(0, 0) \in A$, and for any $\varepsilon > 0$, the open ball $B((0, 0), \varepsilon)$ will contain points with irrational $y$ coordinates, meaning $B((0, 0), \varepsilon) \not\subseteq A$. Alternatively, this set is also countable, so it cannot be the superset of any uncountable open ball. This thing is not closed either: if we take $(\sqrt{2}, \sqrt{2}) \not\in A$ and any open ball $B((\sqrt{2}, \sqrt{2}), \varepsilon)$, we can find a rational point $(q, q)$ in that ball and in $A$. We can also define a sequence like $x_n = (2^{1/n}, 2)$ that approaches $(1, 2) \not\in A$.
\end{enumerate}

\paragraph{Problem 4.}
Let $F \subseteq X$ be any subset of a metric space $(X, d)$. Define \[\rho(x, F) = \inf_{y\in F} d(x, y).\] 
\begin{enumerate}
\item This function is continuous.
\begin{proof}
    Let $U \subseteq \mathbb R$ be any open set, and in its preimage, consider some point $x \in \rho^{-1}(U, F)$. Let $\varepsilon > 0$ be such that we have an open ball $B(\rho(x, F), \varepsilon) \subseteq U$, and using the property of infimums, find some $q \in F$ be such that $d(q, x) < \rho(x, F) + \varepsilon / 2$. 
    
    We wish to show that $\rho(B(x, \varepsilon / 2), F) \subseteq B(\rho(x, F), \varepsilon)$. Let $y \in B(x, \varepsilon / 2)$, and note that by the triangle inequality, \begin{align*} \rho(y, F) &\le d(q, y) \le d(q, x) + d(x, y) \\ &< \rho(x, F) + \varepsilon / 2 + \varepsilon / 2 = \rho(x, F) + \varepsilon.
     \end{align*}
    Similarly, choose some $r \in F$ such that $d(r, y) < \rho(y, F) + \varepsilon / 2$, so
    \begin{align*} 
        \rho(x, F) &\le d(r, x) \le d(r, y) + d(y, x) \\ &< \rho(y, F) + \varepsilon / 2 + \varepsilon / 2 = \rho(y, F) + \varepsilon.
    \end{align*}
    Thus, $\rho(x, F) - \varepsilon < \rho(y, F) < \rho(x, F) + \varepsilon$ for all $y \in B(x, \varepsilon / 2)$. We have shown that $B(x, \varepsilon / 2) \subseteq \rho^{-1}(U, F)$, so the preimage is open and $\rho$ is continuous.
\end{proof}

\item $F$ is closed if and only if $\rho(x, F) > 0$ for all $x \not\in F$. 
\begin{proof}
    ($\implies$) Let $x \not\in F$, and choose some $\varepsilon > 0$ such that the open ball $B(x, \varepsilon)$ does not intersect $F$. By the property of the infimum, choose some $q \in F$ such that $d(q, x) < \rho(x, F) + \varepsilon / 2$, and note that since $q \not\in B(x, \varepsilon)$, we get
    \[\varepsilon \le d(q, x) < \rho(x, F) + \varepsilon / 2, \]
    so $\rho(x, F) > \varepsilon / 2 > 0$ as desired.

    ($\impliedby$) Let $x \not\in F$, and let $\varepsilon$ be such that $0 < \varepsilon < \rho(x, F)$. Consider any point $y \in B(x, \varepsilon)$, and assume for contradiction that $y \in F$. Then, $\rho(x, F) \le d(y, x) < \varepsilon$, contradicting $\varepsilon < \rho(x, F)$. Thus, we get $B(x, \varepsilon) \cap F = \varnothing$.
\end{proof}

\item Let $F_1, F_2 \subseteq X$ be disjoint closed sets. There exist open sets $U_1, U_2 \subseteq X$ such that $U_1 \cap U_2 = \varnothing$ and $F_1 \subseteq U_1$ and $F_2 \subseteq U_2$. (Metric spaces are normal).
\begin{proof}
    Consider the function $\chi : X \to \mathbb R$ defined by $\chi(x) = \rho(x, F_1) - \rho(x, F_2)$. By the continuity of $\rho$, $\chi$ is continuous too. Thus, we can take $U_1 = \chi^{-1}((-\infty, 0))$ and $U_2 = \chi^{-1}((0, -\infty))$. These sets are disjoint open sets as they are the continuous preimages of disjoint open sets. Note that for any $x \in F_1$, we get that $x \not\in F_2$ so $\rho(x, F_1) = 0$ while $\rho(x, F_2) > 0$ by part 2 above. Thus, $\chi(x) < 0$, so $x \in U_1$ and $F_1 \subseteq U_1$ as desired. By a symmetric argument, $F_2 \subseteq U_2$. 
\end{proof}

\end{enumerate}

\end{document}